% common/macros.tex — modo final, cajas de borrador y helpers.
% Depende de: etoolbox, tcolorbox, xcolor (cargados en preamble.tex).

% ===== Modo final =====
% Dos vías equivalentes:
%   1. \documentclass[final]{cysec}
%   2. `make ws-01 FINAL=1` → crea .final-mode en la raíz del proyecto.
% Misma mecánica que Tesis-MACC/project/common/macros.tex.
\newif\iffinalmode
\finalmodefalse
\ifcysec@final\finalmodetrue\fi
\IfFileExists{.final-mode}{\finalmodetrue}{}

% ===== \porque{título}{texto} — nota de borrador, invisible en modo final =====
\newtcolorbox{porquebox}[1]{%
  enhanced, breakable, sharp corners,
  colback=yellow!4, colframe=orange!75!black,
  boxrule=0.4pt, left=6pt, right=6pt, top=4pt, bottom=4pt,
  fonttitle=\bfseries\footnotesize,
  coltitle=white,
  attach boxed title to top left={xshift=4pt,yshift=-2pt},
  boxed title style={colback=orange!75!black, sharp corners, boxrule=0pt,
    top=1pt,bottom=1pt,left=4pt,right=4pt},
  title={\cysec@lbl{Nota de borrador}{Draft note}\, #1}%
}
\newcommand{\porque}[2]{%
  \iffinalmode\else
    \par\smallskip
    \begin{porquebox}{#1}\small #2\end{porquebox}
    \par\smallskip
  \fi
}

% ===== \todoans — marca de respuesta pendiente =====
% Visible en borrador; en modo final desaparece pero deja warning en el .log.
\newcommand{\todoans}[1][]{%
  \PackageWarning{cysec}{Respuesta pendiente\ifstrempty{#1}{}{ (#1)} en la
    pregunta \thepregunta}%
  \iffinalmode\else
    \par\smallskip
    \begin{tcolorbox}[enhanced, sharp corners, colback=yellow!8,
      colframe=orange!80!black, boxrule=0.4pt, left=6pt, right=6pt,
      top=3pt, bottom=3pt]%
      \small\bfseries\color{orange!60!black}%
      \cysec@lbl{Pendiente por responder}{Answer pending}%
      \ifstrempty{#1}{}{\normalfont\ — #1}%
    \end{tcolorbox}%
    \par\smallskip
  \fi
}

% ===== Helpers de texto =====
\newcommand{\term}[1]{\textbf{#1}}                       % concepto clave
\newcommand{\eng}[1]{\foreignlanguage{english}{\emph{#1}}} % término en inglés
